%%
%% This is file `gabarit-these-articles.tex',
%% generated with the docstrip utility.
%%
%% The original source files were:
%%
%% hecthese.dtx  (with options: `gabarit,phd,articles,francais')
%% 
%% This is a stripped version of the original file.
%% 
%% Copyright 2017-2021 HEC Montreal
%% 
%% This work may be distributed and/or modified under the
%% conditions of the LaTeX Project Public License, either version 1.3c
%% of this license or (at your option) any later version.
%% 
%% The latest version of this license is in
%% http://www.latex-project.org/lppl.txt
%% and version 1.3c or later is part of all distributions of LaTeX
%% version 2008/05/04 or later.
%% 
%% This work has the LPPL maintenance status `maintained'.
%% 
%% The Current Maintainer of this work is Benoit Hamel
%% <benoit.2.hamel@hec.ca>.
%% 
%% This work consists of the files hecthese.dtx, hecthese-fr.ins,
%% hecthese-en.ins, hecthese.pdf, hecthese-en.pdf
%% and the derived files listed in the README file.
%% 
%% GABARIT POUR UNE THESE PAR ARTICLES
%%
%% Ceci est le fichier maitre dans lequel vous inscrivez les m\'{e}tadonn\'{e}es
%% relatives \`{a} votre travail, vous cr\'{e}ez vos commandes et environnements
%% personnalis\'{e}s, et \`{a} partir duquel vous lancez vos compilations.
%%
%% NE REDIGEZ PAS VOTRE THESE OU MEMOIRE DANS CE FICHIER!
%%
%% Consultez la documentation de la classe hecthese pour de plus amples
%% informations.
%%
%% DECLARATION DE LA CLASSE DE DOCUMENT
%%
%% La classe est declaree avec le type de document et les langues par
%% defaut. Inscrivez dans la liste d'options la taille de police de
%% caracteres (10pt, 11pt, 12pt) ou laissez la classe charger la taille
%% par defaut : 12pt.
\documentclass[phdarticles,english,french]{hecthese}
%%
%% PACKAGES DEJA CHARGES
%%
%% Voici la liste des packages qui viennent deja charges avec la classe
%% hecthese. VOUS NE DEVEZ PAS LES RECHARGER A NOUVEAU! Si vous le faites,
%% vous occasionnerez des bogues.
%%
%% Package[utf8]{inputenc}
%% Package[T1]{fontenc}
%% Package{natbib}
%% Package{chapterbib}
%% Package{babel}
%% Package[autolanguage]{numprint}
%% Package{calc}
%% Package{enumitem}
%% Package{tocvsec2}
%% Package{graphicx}
%% Package{color}
%% Package{amsmath}
%% Package{iflang}
%%
%% PACKAGES A CHARGER
%%
%% Ajoutez tous les packages necessaires a la redaction de votre travail.
%% Assurez-vous cependant de suivre les consignes suivantes :
%%
%% 1) Le package hyperref doit etre charge EN DERNIER si vous voulez qu'il
%% fonctionne correctement.
%% 2) Le package geometry est INCOMPATIBLE avec la classe memoir. Vous ne
%% devez pas l'utiliser dans votre travail. Consultez la documentation
%% de la classe memoir pour de plus amples informations.
%%
%% CHOIX D'UNE POLICE DE CARACTERES
%%
%% Choisissez le package mathptmx si vous voulez utiliser une police de type
%% Times, avec empattements, et le package mathpazo si vous voulez utiliser
%% une police de type Arial, sans empattements. Choisissez-en une et supprimez
%% l'autre, ou mettez-la en commentaires.
\usepackage{mathptmx}
%% \usepackage{mathpazo}

\usepackage{hyperref}
%%
%% PRODUCTION DE L'INDEX
%%
\makeindex
%%
%% GENERATION DES TITRES
%%
%% On change les titres de l'introduction et de la conclusion generales.
%%
\HECgenererTitres
%%
%%
%% MISE EN FORME DE LA TABLE DES MATIERES
%%
%% On inclut dans la table des matieres toutes les divisions de document
%% jusqu'aux sous-sections. Si vous desirez avoir une table des matieres
%% plus detaillee, indiquez dans les deux commandes ci-dessous jusqu'a
%% quel niveau vous voulez voir repertories dans la TDM.
%%
\setsecnumdepth{subsection} % Numerotation des sous-sections / Subsection numbering
\settocdepth{subsection} % Inclusion des sous-sections dans la TDM / Including subsections in the TOC
%%
%% METADONNEES DU DOCUMENT
%%
%% Le titre de votre travail. Si le titre est long, utilisez la commande \\
%% pour le mettre sur plusieurs lignes.
\HECtitre{Titre de la th\`{e}se ou du m\'{e}moire}
%% Le sous-titre de votre travail. S'il ne comporte pas de sous-titre, videz
%% le contenu des accolades.
\HECsoustitre{Sous-titre de la th\`{e}se}
%% L'auteur, c'est vous...
\HECauteur{Pr\'{e}nom Nom}
%% Nom de l'option de votre grade
\HECoption{Nom de l'option}
%% Mois du depot final de votre travail
\HECmoisDepot{Mai}
%% Annee du depot final de votre travail
\HECanneeDepot{2017}
%% Le nom complet du president rapporteur et son genre (M ou F)
\HECpresidentRapporteur{Pr\'{e}nom Nom}{M ou F}
%% Le nom complet de votre directeur de recherche et son genre (M ou F)
\HECdirecteurRecherche{Pr\'{e}nom Nom}{M ou F}
%% Le nom complet de votre codirecteur de recherche et son genre (M ou F)
\HECcodirecteurRecherche{Pr\'{e}nom Nom}{M ou F}
%% L'universite de provenance de votre codirecteur de recherche
\HECuniversiteCodirecteur{HEC Montr\'{e}al}
%% Le nom complet du membre du jury
\HECmembreJury{Pr\'{e}nom Nom}
%% L'universite de provenance du membre du jury
\HECuniversiteMembreJury{HEC Montr\'{e}al}
%% Le nom complet de l'examinateur externe et son genre (M ou F)
\HECexaminateurExterne{Pr\'{e}nom Nom}{M ou F}
%% L'universite de provenance de l'examinateur externe
\HECuniversiteExaminateur{HEC Montr\'{e}al}
%% Le nom complet du representant du directeur et son genre (M ou F)
\HECrepresentantDirecteur{Pr\'{e}nom Nom}{M ou F}
%%
%% OPTIONS DES PACKAGES CHARGES
%%
%% Si vos packages ont des options specifiques a charger avant le debut
%% du document, inscrivez-les ci-dessous. Si vous voulez outrepasser les
%% options des packages charges par defaut avec la classe hecthese,
%% consultez la documentation pour en connaitre la procedure.
%%
%% Options du package hyperref (inclure les metadonnees pdf dans les options) /
%% hyperref package option (including pdf metadata)
\hypersetup{%
colorlinks=true,
allcolors=black,
pdfauthor=\HECpdfauteur,
pdftitle=\HECpdftitre
}
%% Options de babel / babel options
\frenchbsetup{% 
og=«, fg=» % caracteres « et » sont les guillemets
}
%%
%% DEBUT DE LA THESE OU DU MEMOIRE
%%
\begin{document}
%%
%% STYLE BIBLIOGRAPHIQUE
%%
%% On utilise par defaut le style bibliographique francais issu du package
%% francais-bst. L'utilisation de ce style n'est pas obligatoire. Consultez
%% la documentation pour connaitre la liste des styles compatibles avec la
%% langue de redaction de votre these ou memoire.
%%
\bibliographystyle{francais}

%% Pages liminaires / frontmatter
\frontmatter

%% Page de garde / cover page
\mbox{}
\thispagestyle{empty}
\cleardoublepage

%% Page de titre
 \HECpagestitre

%% Configuration des bibliographies des articles /
%% Articles' bibliographies configuration
\HECbibliographieArticle

%% Resume francais / french abstract
 %% Fichier contenant le resume francais, les mots-cles et les methodes de recherche.
\chapter*{R\'{e}sum\'{e}}
\phantomsection\addcontentsline{toc}{chapter}{R\'{e}sum\'{e}}
\thispagestyle{empty} % Premiere page non paginee / Unnumbered first page

%% Redigez votre resume francais ici (350 a 500 mots).

\section*{Mots-cl\'{e}s}

%% Inscrivez vos mots-cles ici (15 maximum, incluant les methodes de recherche ci-dessous).

\section*{M\'{e}thodes de recherche}

%% Inscrivez vos methodes de recherche ici.

%% THESES ET MEMOIRES PAR ARTICLES SEULEMENT
%% Si vous avez insere des citations dans cette section, retirez les signes
%% de commentaires (%%) devant la commande ci-dessous et inscrivez le style
%% bibliographique et le nom du fichier .bib utilises pour vos references.
%% \HECreferences{style}{nom-du-fichier-file-name}


%% Resume anglais / english abstract
 %% Fichier contenant le resume anglais, les mots-cles et les methodes de recherche.
\chapter*{Abstract}
\phantomsection\addcontentsline{toc}{chapter}{Abstract}

%% Redigez votre resume anglais ici (350 a 500 mots).

\section*{Keywords}

%% Inscrivez vos mots-cles en anglais ici (15 maximum, incluant les methodes de recherche ci-dessous).

\section*{Research methods}

%% Inscrivez vos methodes de recherche en anglais ici.

%% THESES ET MEMOIRES PAR ARTICLES SEULEMENT
%% Si vous avez insere des citations dans cette section, retirez les signes
%% de commentaires (%%) devant la commande ci-dessous et inscrivez le style
%% bibliographique et le nom du fichier .bib utilises pour vos references.
%% \HECreferences{style}{nom-du-fichier-file-name}


%% Table des matieres (* pour ne pas inclure la TDM dans la TDM) /
%% Table of contents (* for excluding the TOC from the TOC)
\tableofcontents*
\cleardoublepage

%% Liste des tableaux / list of table
\listoftables
\cleardoublepage

%% Liste des figures / list of figures
\listoffigures
\cleardoublepage

%% Liste des abreviations / acronyms list
 %% Fichier contenant la liste des abreviations. Vous retrouverez ci-dessous un exemple
%% de liste. L'environnement HECabreviations prend en argument la plus longue des
%% abreviations. A la compilation, une liste en deux colonnes alignees est generee.
\chapter*{\HECtdmAbreviations}
\phantomsection\addcontentsline{toc}{chapter}{\HECtdmAbreviations}

\begin{HECabreviations}{ABBR}
\item[ABBR] Abr\'{e}viation
\item[BAA] Baccalaur\'{e}at en administration des affaires
\item[DESS] Dipl\^{o}me d'\'{e}tudes sup\'{e}rieures sp\'{e}cialis\'{e}es
\item[HEC] Hautes \'{e}tudes commerciales
\item[MBA] Ma\^{i}trise en administration des affaires
\item[MSc] Ma\^{i}trise
\item[PhD] Doctorat
\end{HECabreviations}


%% Dedicace / dedication
 %% Fichier contenant la dedicace
%% N'inscrivez rien entre les accolades de la commande \chapter*{},
%% sauf si vous voulez voir la dedicace dans la table des matieres.

\chapter*{}

\begin{HECdedicace}
%% Redigez votre dedicace ici.
\end{HECdedicace}


%% Remerciements / acknowledgements
 %% Fichier contenant les remerciements
\chapter*{\HECtdmRemerciements}
\phantomsection\addcontentsline{toc}{chapter}{\HECtdmRemerciements}

%% Redigez vos remerciements ici.


%% Avant-propos / preface
 %% Fichier contenant l'avant-propos
\chapter*{\HECtdmAvantPropos}
\phantomsection\addcontentsline{toc}{chapter}{\HECtdmAvantPropos}

%% Redigez votre avant-propos ici.

%% THESES ET MEMOIRES PAR ARTICLES SEULEMENT
%% Si vous avez insere des citations dans cette section, retirez les signes
%% de commentaires (%%) devant la commande ci-dessous et inscrivez le style
%% bibliographique et le nom du fichier .bib utilises pour vos references.
%% \HECreferences{style}{nom-du-fichier-file-name}


%% Partie principale de la these ou du memoire / mainmatter
\mainmatter

%% Introduction
%% Fichier contenant l'introduction
\chapter*{\HECtitreIntroduction}
\phantomsection\addcontentsline{toc}{chapter}{\HECtitreIntroduction}
\thispagestyle{empty} % Premiere page non paginee / First page is unnumbered

%% Redigez votre introduction ici.

%% THESES ET MEMOIRES PAR ARTICLES SEULEMENT
%% Si vous avez insere des citations dans cette section, retirez les signes
%% de commentaires (%%) devant la commande ci-dessous et inscrivez le style
%% bibliographique et le nom du fichier .bib utilises pour vos references.
%% \HECreferences{style}{nom-du-fichier}


%% Cadre theorique / theoretical framework
 %% Fichier contenant le cadre theorique
\chapter*{\HECtdmCadreTheorique}
\phantomsection\addcontentsline{toc}{chapter}{\HECtdmCadreTheorique}
\thispagestyle{empty}

%% Redigez votre cadre theorique ici.

%% THESES ET MEMOIRES PAR ARTICLES SEULEMENT
%% Si vous avez insere des citations dans cette section, retirez les signes
%% de commentaires (%%) devant la commande ci-dessous et inscrivez le style
%% bibliographique et le nom du fichier .bib utilises pour vos references.
%% \HECreferences{style}{nom-du-fichier-file-name}


%% Articles de developpement / articles
%% Fichier contenant un article. La classe genere trois fichiers
%% d'articles par defaut. Une these compte generalement trois articles
%% et un memoire, un. Si vous en avez besoin davantage,
%% enregistrez ce fichier sous un autre nom et incluez-le dans
%% votre gabarit avec la commande \include.
%%
%% Les articles sont structures tel que vous le voyez ci-dessous :
%% un resume non numerote, une introduction, des sections et une
%% conclusion numerotees, et une bibliographie. Vous pouvez ajouter
%% ou supprimer des sections de developpement selon vos besoins.
\chapter{Titre de l'article / Article title}
\thispagestyle{empty} % Premiere page non paginee / First page is unnumbered

\section*{\HECtdmResumeArticle}
\phantomsection\addcontentsline{toc}{section}{\HECtdmResumeArticle}

%% Redigez votre resume ici.

\section{Introduction}

%% Redigez votre introduction d'article ici.

\section{Titre de la section de d\'{e}veloppement 1 / Section 1 title}

%% Redigez votre section de developpement ici.

\section{Titre de la section de d\'{e}veloppement 2 / Section 2 title}

%% Redigez votre section de developpement ici.

\section{Titre de la section de d\'{e}veloppement 3 / Section 3 title}

%% Redigez votre section de developpement ici.

\section{Conclusion}

%% Redigez votre conclusion d'article ici.

\bibliographystyle{francais}
%% Inscrivez le nom de votre fichier .bib entre les accolades.
\bibliography{}

%% Fichier contenant un article. La classe genere trois fichiers
%% d'articles par defaut. Une these compte generalement trois articles
%% et un memoire, un. Si vous en avez besoin davantage,
%% enregistrez ce fichier sous un autre nom et incluez-le dans
%% votre gabarit avec la commande \include.
%%
%% Les articles sont structures tel que vous le voyez ci-dessous :
%% un resume non numerote, une introduction, des sections et une
%% conclusion numerotees, et une bibliographie. Vous pouvez ajouter
%% ou supprimer des sections de developpement selon vos besoins.
\chapter{Titre de l'article / Article title}
\thispagestyle{empty} % Premiere page non paginee / First page is unnumbered

\section*{\HECtdmResumeArticle}
\phantomsection\addcontentsline{toc}{section}{\HECtdmResumeArticle}

%% Redigez votre resume ici.

\section{Introduction}

%% Redigez votre introduction d'article ici.

\section{Titre de la section de d\'{e}veloppement 1 / Section 1 title}

%% Redigez votre section de developpement ici.

\section{Titre de la section de d\'{e}veloppement 2 / Section 2 title}

%% Redigez votre section de developpement ici.

\section{Titre de la section de d\'{e}veloppement 3 / Section 3 title}

%% Redigez votre section de developpement ici.

\section{Conclusion}

%% Redigez votre conclusion d'article ici.

\bibliographystyle{francais}
%% Inscrivez le nom de votre fichier .bib entre les accolades.
\bibliography{}

%% Fichier contenant un article. La classe genere trois fichiers
%% d'articles par defaut. Une these compte generalement trois articles
%% et un memoire, un. Si vous en avez besoin davantage,
%% enregistrez ce fichier sous un autre nom et incluez-le dans
%% votre gabarit avec la commande \include.
%%
%% Les articles sont structures tel que vous le voyez ci-dessous :
%% un resume non numerote, une introduction, des sections et une
%% conclusion numerotees, et une bibliographie. Vous pouvez ajouter
%% ou supprimer des sections de developpement selon vos besoins.
\chapter{Titre de l'article / Article title}
\thispagestyle{empty} % Premiere page non paginee / First page is unnumbered

\section*{\HECtdmResumeArticle}
\phantomsection\addcontentsline{toc}{section}{\HECtdmResumeArticle}

%% Redigez votre resume ici.

\section{Introduction}

%% Redigez votre introduction d'article ici.

\section{Titre de la section de d\'{e}veloppement 1 / Section 1 title}

%% Redigez votre section de developpement ici.

\section{Titre de la section de d\'{e}veloppement 2 / Section 2 title}

%% Redigez votre section de developpement ici.

\section{Titre de la section de d\'{e}veloppement 3 / Section 3 title}

%% Redigez votre section de developpement ici.

\section{Conclusion}

%% Redigez votre conclusion d'article ici.

\bibliographystyle{francais}
%% Inscrivez le nom de votre fichier .bib entre les accolades.
\bibliography{}


%% Conclusion
%% Fichier contenant la conclusion
\chapter*{\HECtitreConclusion}
\phantomsection\addcontentsline{toc}{chapter}{\HECtitreConclusion}
\thispagestyle{empty} % Premiere page non paginee / First page is unnumbered

%% Redigez votre conclusion ici.

%% THESES ET MEMOIRES PAR ARTICLES SEULEMENT
%% Si vous avez insere des citations dans cette section, retirez les signes
%% de commentaires (%%) devant la commande ci-dessous et inscrivez le style
%% bibliographique et le nom du fichier .bib utilises pour vos references.
%% \HECreferences{style}{nom-du-fichier-file-name}


%% Index analytique / analytical index
\printindex

%% BIBLIOGRAPHIE / BIBLIOGRAPHY
 %% Configuration de la bibliographie generale
\HECbibliographieGenerale
%% Inscrivez le nom de votre fichier .bib entre les accolades.
\bibliography{}

\backmatter

%% Retour a la pagination romaine / Back to roman page numbering
\pagenumbering{roman}

%% Annexes / appendices
\appendix
 %% Fichier contenant une annexe. La classe ne genere qu'un seul
%% fichier annexe.tex par defaut. Si vous en avez besoin davantage,
%% enregistrez ce fichier sous un autre nom et incluez-le dans votre
%% gabarit avec la commande \include.
%%
%% Pour que la bibliographie et la (les) annexe(s) soient paginees
%% correctement, c'est-a-dire en chiffres arabes pour la bibliographie
%% et en chiffres romains pour les annexes, ces dernieres doivent etre
%% placees apres la commande \backmatter. Ce faisant, la numerotation
%% des annexes s'en trouve desactivee. Vous devrez donc numeroter
%% vos annexes manuellement a l'interieur de la commande \chapter.
\chapter{Annexe A -- Titre de l'annexe}

%% Redigez votre annexe ici.

%% THESES ET MEMOIRES PAR ARTICLES SEULEMENT
%% Si vous avez insere des citations dans cette section, retirez les signes
%% de commentaires (%%) devant la commande ci-dessous et inscrivez le style
%% bibliographique et le nom du fichier .bib utilises pour vos references.
%% \HECreferences{style}{nom-du-fichier-file-name}


%% Page de garde de fin / back cover page
\mbox{}
\thispagestyle{empty}

\end{document}
\endinput
%%
%% End of file `gabarit-these-articles.tex'.
